\section{Related Work}
\label{sec:related}

\paragraph{Financial forecasting from price and text.}
Sequence models over OHLCV data, including Transformer-based encoders~\cite{vaswani2017attention,zhou2021informer,wu2021autoformer}, temporal fusion models~\cite{lim2021tft}, and newer long-horizon backbones such as PatchTST, TimesNet, and iTransformer~\cite{nie2023patchtst,wu2023timesnet,liu2024itransformer}, remain strong baselines for short-horizon financial prediction, but they are structurally limited to endogenous market dynamics.
Text-augmented models extend this setting by incorporating exogenous signals from news, social media, or earnings calls.
StockNet~\cite{xu2018stocknet} jointly models price histories and Twitter text, while event-driven approaches such as Ding et al.~\cite{ding2015deepevent} show that structured news events can improve prediction.
More recent financial language models, including FinBERT~\cite{yang2020finbert}, FinGPT~\cite{lopez2023fingpt}, and BloombergGPT~\cite{wu2023bloomberggpt}, further improve text-side representation quality.
However, most of this literature assumes that text and price can be synchronously paired, either by aggregating text into fixed periods or by restricting to narrow publication windows.
That assumption is weak for cryptocurrency markets, where web intelligence arrives continuously and with highly variable latency.

\paragraph{Multimodal fusion and the asynchronous-evaluation gap.}
While the financial forecasting literature provides strong unimodal and text-augmented baselines, it largely sidesteps the asynchronous challenge that broader multimodal fusion research has only recently begun to surface.
More broadly, multimodal learning has been organized around representation, translation, alignment, fusion, and co-learning problems~\cite{baltrusaitis2019multimodal}, with gating-based and memory-based fusion families such as GMU, MFN, MISA, and DMNC~\cite{arevalo2017gmu,zadeh2018mfn,hazarika2020misa,le2018dmnc} highlighting the importance of selective cross-modal interaction over time and MulT~\cite{tsai2019multimodal} making the unaligned-sequence setting especially visible for language-centered benchmarks.
Multimodal financial systems such as MDRM~\cite{qin2019mdrm} and the earnings-based risk framework of Sawhney et al.~\cite{sawhney2020multimodal} demonstrate that joint modeling across modalities can improve downstream performance, but they operate in settings with tightly scheduled events and much easier synchronization.
General-purpose fusion architectures therefore provide the main methodological reference points for our comparisons: BiLSTM late fusion~\cite{xu2018stocknet}, directional cross-modal attention in MulT~\cite{tsai2019multimodal}, multiplicative interaction modeling in TFN~\cite{zadeh2017tensor}, modality gating in GMU~\cite{arevalo2017gmu}, memory-based sequential fusion in MFN~\cite{zadeh2018mfn}, and the broader fusion taxonomy summarized by Baltru\v{s}aitis et al.~\cite{baltrusaitis2019multimodal}.
In cryptocurrency prediction, prior multimodal datasets and systems remain mostly daily and do not expose within-day freshness.
Sentiment from Twitter~\cite{abraham2018cryptocurrency}, Reddit~\cite{kim2019predicting}, fear\,\&\,greed indicators~\cite{chen2021fear}, and on-chain activity~\cite{kondor2014rich} has been studied separately, while integrated benchmarks such as PreBit~\cite{zou2023prebit}, DAM~\cite{li2024dam}, CryptoPulse~\cite{liu2024cryptopulse}, and Barth{\'e}lemy et al.~\cite{barthelemy2025deepnlp} operate at daily granularity.
As summarized in Table~\ref{tab:benchmark_comparison}, these datasets do not record modality lag $\tau_{\mathrm{lag}}$ per aligned sample, which makes lag-conditioned training and lag-stratified analysis impossible.
\dataset{} is designed to fill this broader asynchronous-evaluation gap by providing 15-minute price alignment together with explicit lag metadata; cryptocurrency markets serve here as a convenient stress-test environment rather than as the paper's only intended downstream domain.

\begin{table*}[t]
\caption{Comparison of cryptocurrency multimodal datasets. \emph{Lag $\tau$} indicates whether per-sample modality lag is recorded; \emph{Rolling+SR} indicates rolling walk-forward evaluation with Sharpe reporting. LLM synthesis trades source authenticity for temporal control and sub-hourly coverage; the associated limitations are discussed in Section~\ref{sec:limitations}.}
\label{tab:benchmark_comparison}
\centering
\small
\setlength{\tabcolsep}{4pt}
\renewcommand{\arraystretch}{1.05}
\begin{tabular}{llccccc}
\toprule
Dataset & Text source & Price granularity & Assets & Samples & Lag $\tau$ recorded & Rolling+SR \\
\midrule
PreBit~\cite{zou2023prebit}            & Real (Twitter)          & Daily   & BTC         & ${\sim}$1K         & \texttimes & \texttimes \\
DAM~\cite{li2024dam}                   & Real (CryptoCompare)    & Daily   & BTC         & ${\sim}$1K         & \texttimes & \texttimes \\
CryptoPulse~\cite{liu2024cryptopulse}  & Real (Cointelegraph)    & Daily   & Top-5       & ${\sim}$1K         & \texttimes & \texttimes \\
Barth{\'e}lemy et al.~\cite{barthelemy2025deepnlp} & Real (Twitter/Reddit) & Daily & BTC/ETH  & ${\sim}$3yr        & \texttimes & \checkmark \\
\midrule
\dataset{} (synthetic)                & LLM-synthesized         & \textbf{15 min} & BTC/ETH/SOL & \textbf{21{,}374} & \checkmark & \checkmark \\
\dataset{} (real news)                & Real (CryptoCompare)    & \textbf{15 min} & BTC/ETH/SOL & \textbf{27{,}914} & \checkmark & \checkmark \\
\bottomrule
\end{tabular}
\end{table*}

\paragraph{Temporal alignment, freshness, and lag-aware trust.}
Beyond finance, related ideas appear in temporal knowledge graph embedding, temporally weighted recommendation, and sequence models that explicitly acknowledge unaligned context windows.
HyTE~\cite{dasgupta2018hyte} conditions representations on validity intervals, temporally weighted co-attention methods in recommendation down-weight stale context~\cite{liu2021temporal}, and MulT-style cross-modal attention shows that alignment need not be hard-coded at equal timestamps~\cite{tsai2019multimodal}.
However, these lines of work stop short of a learned, content-aware trust mechanism that jointly conditions on cross-modal agreement and lag.
What is missing in multimodal financial forecasting is therefore not just a better text encoder, but an explicit mechanism for deciding \emph{when} text should be trusted relative to the current market state.
Existing financial fusion models typically ignore modality lag or treat it as a nuisance variable rather than as an architectural input.
\CGCMA{} takes the opposite view: text first grounds into the price history through cross-modal attention, and lag-aware trust is then handled explicitly by a conditional gate that receives $\tau_{\mathrm{lag}}$ together with the modality representations.
